%% Nature Computational Science Brief Communication draft
%% Initial-submission PDF target; preserve the PRL-style manuscript separately.

\documentclass[aps,prl,twocolumn,superscriptaddress,floatfix]{revtex4-2}

\usepackage{amsmath,amssymb,bm}
\usepackage{graphicx}
\usepackage{xcolor}
\usepackage[colorlinks=true,linkcolor=blue,citecolor=blue,urlcolor=blue]{hyperref}
\usepackage{lineno}

\newcommand{\HAB}{H_{AB}}
\newcommand{\pacc}{p_{\mathrm{acc}}}
\newcommand{\Xnew}{X^{\mathrm{n}}}
\newcommand{\Xold}{X^{\mathrm{o}}}

\begin{document}
\linenumbers

\title{Path Sampling Diagnoses Protein Diffusion Models}

\author{Muhammad R.~Hasyim}
\affiliation{Simons Center for Computational Physical Chemistry,
New York University, New York, NY~10003}
\email{mh7373@nyu.edu}

\date{\today}

\begin{abstract}
Diffusion models can generate protein--ligand structures in seconds, but one
prediction does not reveal the stochastic denoising paths that led to it. I
adapt transition path sampling to explicit-noise structure diffusion and add
on-the-fly probability enhanced sampling to map ligand-pose denoising
landscapes. The method records the random variables used by Boltz-2 reverse
diffusion, assigns exact path weights and resamples complete denoising
trajectories. For a held-out MEK1--FZC complex, the biased trajectory ensemble
contains a dominant native-like pose and persistent secondary basins with
ligand RMSD above 1~\AA. These rare basins do not appear in single forward
passes but remain stable over tens to hundreds of Monte Carlo steps. The
result is a path-level test for memorization, ambiguity and physical
plausibility in generative structure models.
\end{abstract}

\maketitle

Diffusion-based structure models such as AlphaFold~3 and Boltz-2 now produce
protein complexes and ligand poses quickly enough to be used as routine
screening tools~\cite{alphafold3,boltz2}. That speed leaves a diagnostic gap.
When a generated pose looks chemically plausible, it is hard to tell whether
the model sampled a physically meaningful binding mode or returned a memorized
training-set attractor. Recent protein--ligand benchmarks report strong
dependence on similarity to known complexes and adversarial failures in which
co-folding models preserve ligand placement after binding-site mutations
remove the expected interactions~\cite{masters2024,skrinjar2025}. A useful
test should interrogate the distribution of denoising trajectories for a fixed
input, not just the final structure from one random seed.

Transition path sampling (TPS) is built for that kind of question. Classical
TPS samples rare dynamical trajectories connecting two state volumes rather
than configurations drawn independently~\cite{dellago1998,bolhuis2002}. A
reverse diffusion sampler has the same minimal ingredients when its random
variables are explicit. At denoising step $i$, Boltz-2 draws an SE(3)
augmentation and Gaussian churn noise, then maps the current coordinates
deterministically to the next coordinates. Recording these draws gives a path
density that factorizes over known priors. Forward shooting resamples the
suffix of a denoising path; backward shooting proposes a new prefix with a
Hastings correction; global reshuffling draws a fresh trajectory from the
prior. Supplementary Notes 1, 2 and 5 give the derivation, including Jacobian
cancellations, non-reversible proposals and the relation to quotient-space
diffusion.

Plain TPS samples the unbiased reactive ensemble, so it stays near whichever
basin dominates the denoising landscape. To reach rare terminal poses, I
coupled TPS to on-the-fly probability enhanced sampling
(OPES)~\cite{invernizzi2020,invernizzi2022}. OPES builds an adaptive bias on
terminal collective variables, here ligand heavy-atom RMSD $s_1$ and
ligand--pocket distance $s_2$, and penalizes regions the chain has already
visited while preserving a reweightable path ensemble. This produces one
biased trajectory ensemble on a projected denoising landscape, not a set of
unconnected Boltz-2 predictions.

\begin{figure}[t]
\centering
\fbox{\parbox{0.94\columnwidth}{\centering
\vspace{0.45in}
Placeholder for the final three-panel OPES--TPS landscape: OPES kernel
mixture, reweighted density and visited $(s_1,s_2)$ histogram.
\vspace{0.45in}}}
\caption{\textbf{TPS+OPES maps the terminal ligand-pose landscape.}
The two-dimensional collective variable combines ligand pose RMSD $s_1$ and
ligand--pocket distance $s_2$. In the MEK1--FZC case study, the run comprised
11,590 Monte Carlo steps, 3,184 accepted moves and a 27.5\% acceptance rate.
The projected landscape contains a dominant native-like basin near
$(0.07~\text{\AA}, 5.37~\text{\AA})$ and secondary high-RMSD regions
around larger pocket distances. Numerical definitions and OPES reweighting are
reported in Supplementary Notes 3 and 4.}
\label{fig:landscape}
\end{figure}

I applied this protocol to the MEK1 kinase MAP2K1 with ligand FZC from PDB
7XLP, a post-training-cutoff protein--ligand system highlighted in recent
co-folding analyses as lying in the low-similarity regime where memorized
ligand placement is expected to be least reliable~\cite{masters2024}. A
one-dimensional TPS+OPES run using only $s_1$ already showed a broad tail:
4,180 Monte Carlo steps sampled ligand RMSD from near zero to 1.3~\AA, with
approximately 47\% of accepted paths above 0.5~\AA. The terminal denoising
landscape therefore was not a single rigid pose. The RMSD coordinate alone,
however, could not identify how the ligand moved.

The two-dimensional run resolves this ambiguity (Fig.~\ref{fig:landscape}).
With biasfactor $\gamma=10$, 11,590 Monte Carlo steps produced 3,184 accepted
paths and 1,159 OPES depositions compressed to 82 kernels. The dominant basin,
B$_0$, is native-like at $s_1 \approx 0.07$~\AA\ and
$s_2 \approx 5.37$~\AA. Three additional regions are reproducibly visited:
B$_1$ at $s_1 \approx 0.35$~\AA, B$_2$ at
$s_1 \approx 0.75$~\AA\ with the ligand center displaced by about
0.19~\AA, and a far-tail basin B$_3$ near $s_1 \approx 1.10$~\AA. Across all
logged steps, 8.8\% exceeded 1~\AA\ ligand RMSD and these high-RMSD visits
occurred in persistent runs, with mean dwell length about 73 Monte Carlo steps
and a maximum of 196. The result is not simply stochastic jitter around a
single pose: the model can stabilize alternative terminal geometries.

\begin{figure}[t]
\centering
\fbox{\parbox{0.94\columnwidth}{\centering
\vspace{0.45in}
Placeholder for basin-resolved structural ensemble: B$_0$, B$_1$, B$_2$ and
B$_3$ terminal structures with ligand sticks and protein cartoon/surface.
\vspace{0.45in}}}
\caption{\textbf{Representative structures from denoising basins.}
B$_0$ is a tightly packed native-like pose, B$_1$ shows partial ligand slippage,
B$_2$ is a displaced pocket-engagement mode and B$_3$ contains partially
egressed ligand configurations with subtle pocket-rim loop deformation.
Selection criteria and metric definitions are given in Supplementary Note 3.}
\label{fig:basins}
\end{figure}

Structural inspection of representative terminal frames supports this
interpretation (Fig.~\ref{fig:basins}). B$_0$ retains the compact binding pose,
B$_1$ slips along the pocket axis with the protein fold intact, B$_2$ changes
side-chain contact patterns, and B$_3$ partially exposes the ligand at the
pocket rim. These configurations should not be interpreted immediately as
physical alternative binding modes. The free-energy surface is projected onto
two collective variables, OPES convergence is finite, and independent energy or
geometry checks remain necessary. The point is narrower: a model may look
highly focused under ordinary sampling while still containing rare, stable
denoising basins that matter for uncertainty, memorization and downstream
drug-discovery decisions.

Standard confidence scores and direct multi-seed sampling still matter, but
they do not ask which denoising paths connect noise to structure. The TPS
formulation does. It applies whenever a generator exposes its stochastic
inputs and their densities, including score-based diffusion and related
generative samplers. The framework can compare models, MSA depths or ligand
scaffolds by measuring how the reactive trajectory ensemble changes, rather
than only whether the top-ranked pose matches a reference. The 228-residue
co-folding control, additional diagnostics, evaluation metrics and
computational workflow are reported in the Supplementary Information.

\section*{Methods}

\paragraph{System and sampler.}
The main case study used Boltz-2 in protein--ligand co-folding mode for MEK1
(MAP2K1) with ligand FZC from PDB 7XLP. Reverse diffusion was treated as a
fixed-length Markov chain over coordinates and recorded random variables. State
A is the high-noise start of denoising and state B is the terminal structured
Boltz-2 output. Software, command-line workflow and output files are listed in
Supplementary Note 4.

\paragraph{TPS and OPES settings.}
TPS moves combined forward shooting, backward shooting and global reshuffling.
Forward shooting has unit Metropolis ratio after the endpoint test because the
resampled suffix is drawn from the same prior as the target path density.
Backward shooting uses the non-reversible Hastings correction derived in
Supplementary Note 1. OPES was applied to terminal-frame CVs with biasfactor
$\gamma=10$ and barrier $5\,k_BT$; the two-dimensional run used 11,590 Monte
Carlo steps.

\paragraph{Metrics and availability.}
The primary CVs were ligand heavy-atom RMSD to the reference pose and
ligand-center-to-pocket C$\alpha$ distance. Additional diagnostics include
contact order, steric clashes, lDDT, Ramachandran outliers, radius of gyration,
PoseBusters-style ligand checks and ensemble-distribution metrics
(Supplementary Note 3). Code is available at
\url{https://github.com/muhammadhasyim/tps-diffusion-model}. Large language
model assistance was used for manuscript restructuring and prose drafting; the
author reviewed all scientific content and remains accountable for the work.

\begin{acknowledgments}
The author acknowledges support from the Simons Center for Computational
Physical Chemistry at New York University and NYU IT High Performance
Computing resources.
\end{acknowledgments}

\begin{thebibliography}{25}

\bibitem{alphafold3}
J.~Abramson \textit{et al.}, \textit{Nature} \textbf{630}, 493 (2024).

\bibitem{boltz2}
J.~Wohlwend, G.~Corso, S.~Passaro, M.~Reveiz, V.~Portnoi, R.~Gondal,
T.~Hothorn, R.~Barzilay and T.~Jaakkola, \textit{bioRxiv}
(2025), \url{https://doi.org/10.1101/2025.05.06.651165}.

\bibitem{masters2024}
S.~L.~Masters \textit{et al.}, \textit{bioRxiv}
(2024), \url{https://doi.org/10.1101/2024.06.03.597219}.

\bibitem{skrinjar2025}
P.~\v{S}krinjar, J.~Eberhardt, G.~Tauriello, T.~Schwede and J.~Durairaj,
\textit{bioRxiv} (2025), \url{https://doi.org/10.1101/2025.02.03.636309}.

\bibitem{dellago1998}
C.~Dellago, P.~G.~Bolhuis, F.~S.~Csajka and D.~Chandler,
\textit{J.~Chem.~Phys.} \textbf{108}, 1964 (1998).

\bibitem{bolhuis2002}
P.~G.~Bolhuis, D.~Chandler, C.~Dellago and P.~L.~Geissler,
\textit{Annu.~Rev.~Phys.~Chem.} \textbf{53}, 291 (2002).

\bibitem{karras2022}
T.~Karras, M.~Aittala, T.~Aila and S.~Laine,
\textit{Adv.~Neural Inf.~Process.~Syst.} \textbf{35}, 26565 (2022).

\bibitem{invernizzi2020}
M.~Invernizzi and M.~Parrinello,
\textit{J.~Phys.~Chem.~Lett.} \textbf{11}, 2731 (2020).

\bibitem{invernizzi2022}
M.~Invernizzi and M.~Parrinello,
\textit{J.~Chem.~Theory Comput.} \textbf{18}, 3988 (2022).

\bibitem{song2021}
Y.~Song, J.~Sohl-Dickstein, D.~P.~Kingma, A.~Kumar, S.~Ermon and B.~Poole,
\textit{ICLR} (2021).

\bibitem{xu2026}
Y.~Xu, Y.~Wang, S.~Luo, K.~Gao, T.~He, D.~He and C.~Liu,
\textit{ICLR} (2026), arXiv:2604.21809.

\bibitem{ops}
D.~W.~Swenson, J.-H.~Prinz, F.~Noe, J.~D.~Chodera and P.~G.~Bolhuis,
\textit{J.~Chem.~Theory Comput.} \textbf{15}, 813 (2019).

\bibitem{mariani2013lddt}
V.~Mariani, M.~Biasini, A.~Barbato and T.~Schwede,
\textit{Bioinformatics} \textbf{29}, 2722 (2013).

\bibitem{posebusters2024}
C.~Bouysset and S.~Fiorucci, \textit{J.~Cheminform.} \textbf{13}, 72 (2021).

\end{thebibliography}

\end{document}
